\chapter{The Role and Cost of Preconditioning}\label{v0ch:16}
\begin{chaptersummary}
While Krylov methods build corrections solely from operator applications, preconditioning reshapes directional scales and global coupling. Left, right, and symmetric formulations connect approximate factorizations to subspace projection, and their practical value depends on convergence speedup, inverse-action cost, and amortization across solves, with conditioning, eigenvalue clustering, and nonnormality playing distinct roles across matrix classes.
\end{chaptersummary}
\learninggoals{Distinguish left, right, and symmetric preconditioning; apply inverse actions correctly; include setup and application in total cost.}
\paragraph{Core Question}How can a cheaper problem improve the original solve without changing its exact solution?

\section{Transformations and Residuals}
A fixed invertible preconditioner $M$ enters as
\[
 M^{-1}A\bm x=M^{-1}\bm b\quad\text{(left)},\qquad
 AM^{-1}\bm y=\bm b,\ \bm x=M^{-1}\bm y\quad\text{(right)}.
\]
When $M=CC^T$ is SPD, symmetric preconditioning solves
$C^{-1}AC^{-T}\widehat{\bm x}=C^{-1}\bm b$ followed by $\bm x=C^{-T}\widehat{\bm x}$; in all cases, triangular or approximate solves are applied without forming $M^{-1}$ explicitly.
For SPD systems, spectral equivalence
$c\,\bm v^TM\bm v\le\bm v^TA\bm v\le C\,\bm v^TM\bm v$
bounds the preconditioned condition number by $C/c$, whereas for nonsymmetric operators eigenvector conditioning and nonnormality also govern GMRES residual polynomials.
\conceptbridge{Why can left and right preconditioning have identical eigenvalues but different residual histories?}{For invertible $M$, the left- and right-preconditioned operators are similar and therefore share eigenvalues, yet their solution variables, residual norms, and starting vectors differ because similarity transformations do not preserve Euclidean norms. Specifying both the minimized norm and the original-system acceptance criterion prevents ambiguity.}{Equivalent exact solution $\longrightarrow$ different projection geometries $\longrightarrow$ different residual histories.}
\section{A Minimal Example and Cost}
For $A=\diag(1,10^6)$, choosing $M=A$ eliminates the scale disparity in a single step, whereas for a general large sparse system applying $M=A$ is as hard as the original solve. In practice, Jacobi captures diagonal scaling, block methods capture local coupling, and multigrid addresses cross-scale coupling, with total runtime structured as
\[
 T_{\rm total}=T_{\rm setup}
 +k(T_A+T_{M^{-1}}+T_{\rm orth}+T_{\rm reduce})
 +T_{\rm verify}.
\]
Solving multiple right-hand sides or consecutive time steps amortizes setup cost, whereas significant coefficient changes prompt re-evaluating preconditioner reuse.
\costnote{Effective preconditioning balances memory, setup overhead, application cost, parallel communication, and robustness; a method with fewer iterations but costly sequential triangular solves can take longer than more iterations of a cheaper, more parallel preconditioner.}
\aiexercise{Implement left- and right-preconditioned GMRES for a nonsymmetric three-by-three system. Plot true and preconditioned residuals and explain the difference.}{Given setup and per-step costs for two preconditioners, find the number of right-hand sides at which the more expensive setup pays off.}
\studyroute{Field-of-values bounds and general spectral-equivalence proofs are omitted. See \cite{saad,templates}.}{left preconditioning; right preconditioning; spectral equivalence; time to solution}
